problem:  
**Problem (Deformation–Rigidity of the Heisenberg Groupoid).**  
Let \( M \) be a compact smooth manifold.  
For each \( t\in[0,1] \), let \( \mathcal H_t\subset TM \) be a smooth rank‑\(r\) distribution, endowed with a smooth inner product.  
Suppose \( \mathcal H_t \) satisfies Hörmander’s bracket‑generating condition for all \(t\) outside a finite set \( \{t_1,\dots,t_k\} \), and that the nilpotent graded type of the osculating Lie algebra \( \mathfrak g(\mathcal H_t)_x \) at each point \(x\in M\) is **piecewise constant** in \( t \), changing only at these singular values.

For each regular \( t\), write \(T_{\mathcal H_t}M\) for the osculating groupoid; its convolution \(C^*\)-algebra \(A_t := C^*(T_{\mathcal H_t}M)\) carries a canonical Heisenberg \(K\)-class 
\[
[\sigma_H(P_t)] \in K_0(A_t)
\]
associated to any Rockland operator \(P_t\) compatible with \( \mathcal H_t\).

**Problem:**  
Construct, if possible, a *universal deformation groupoid* \( \mathbb T \rightrightarrows M\times[0,1] \) whose restriction to each fiber \(M\times\{t\}\) recovers \(T_{\mathcal H_t}M\).  
Assuming such \( \mathbb T\) exists, its \(C^*\)-algebra gives a continuous field \(A=\{A_t\}\).  
Investigate whether the section \(t\mapsto [\sigma_H(P_t)]\) defines a continuous class \( [\sigma_H] \in K_0(A)\).  
Is it necessarily true that the induced endpoint classes coincide under the canonical comparisons  
\[
[\sigma_H(P_0)] \;=\; [\sigma_H(P_1)] \quad\text{in}\quad K_0(A_0)\cong K_0(A_1)?
\]
Or can loss of the Hörmander property at an intermediate \(t_j\) change the \(K\)-theory class, yielding a new *Heisenberg rigidity invariant* of filtered‑structure degenerations?

Why it is a "good" problem:  
This formulation cleanly separates existence and invariance issues.  
It is concrete—asking for a deformation groupoid unifying osculating structures and for continuity of the associated \(K\)-class—yet reaching deeply into unsolved territory.  
If such a groupoid exists and continuity fails, that failure would represent a genuine topological obstruction in the noncommutative geometry of hypoelliptic operators.  
If continuity always holds, it would prove a strong rigidity theorem unifying local Hörmander analysis with global \(C^*\)-algebraic topology.  
Either outcome would reshape our understanding of how analytic hypoellipticity interacts with topological invariance, sitting precisely at the intersection of microlocal analysis, sub‑Riemannian geometry, and \(K\)-theory of groupoids.  
This geometric–analytic simplicity hiding profound depth marks it as a **good** research problem.